\section{Methodology}
\label{sec:method}

\Cref{fig:pipeline} summarises the pipeline: selecting post-cutoff pull
requests, injecting defects, admitting variants by execution, retrieving
context, and judging.

\begin{figure}[t]
\centering
\small
\fbox{\parbox{0.95\columnwidth}{%
\textbf{1. Select} PRs created $\geq$ 2025-01-01 (SWE-bench-Live, SWE-rebench)
$\rightarrow$
\textbf{2. Apply} gold patch in the task's Docker image, re-render with \texttt{git diff}
$\rightarrow$
\textbf{3. Inject} one AST mutation per defect class in the region the patch touches
$\rightarrow$
\textbf{4. Admit} a variant only if every test passing with the gold patch still passes
$\rightarrow$
\textbf{5. Retrieve} neutral context from the parent commit
$\rightarrow$
\textbf{6. Judge} (both orders) and analyse.}}
\caption{RepoJudgeBench construction and evaluation pipeline.}
\label{fig:pipeline}
\end{figure}

\subsection{Source pull requests}
We take real, merged pull requests from two continuously updated
repository-level task collections that ship a runnable Docker environment for
every task: the \texttt{full} split of
SWE-bench-Live~\cite{zhang2025swebenchlive} and the \texttt{filtered} and
\texttt{test} splits of SWE-rebench.  Each task provides the gold patch, the
test patch, the issue text, the list of tests that the fix makes pass
(FAIL\_TO\_PASS) and the test command.  To limit training-data contamination we
keep only tasks created on or after 1~January~2025, following the time-window
idea of LiveCodeBench and SWE-bench-Live.  After de-duplication the pool holds
\NPoolInstances{} tasks from \NPoolRepos{} repositories.  We discard gold patches
with more than 400 changed lines and apply a cheap regular-expression screen that
keeps tasks whose diff contains at least one construct targeted by a mutation
operator (e.g.\ an \texttt{except} clause, a \texttt{with} statement, a lock or
an ORM call), in the text the patch adds or in its context lines.

\subsection{Defect taxonomy and injection operators}
\Cref{tab:operators} lists the six defect classes of the proposal and their
operators.  Each class is chosen because the defect is typically invisible to
unit tests but flagged by a careful reviewer.  Operators work on the Python AST
of the \emph{post-gold} file and edit only source lines, so the variant is
syntactically valid and idiomatic.  A site is eligible only if it lies inside
the region the gold patch touches: a line added by the patch, or the body of a
function the patch modifies.  We record whether the mutated site was written by
the patch (\emph{added}) or is pre-existing code in a modified function
(\emph{context}).  Candidates whose AST is unchanged are rejected.  For each
task and class we test up to four candidates (one per operator first, preferring
added sites, in a seeded random order) and admit at most two variants per class,
produced by different operators; at most three classes are tried per task.

\begin{table}[t]
\caption{Defect classes and injection operators.}
\label{tab:operators}
\centering
\footnotesize
\begin{tabular}{@{}>{\raggedright\arraybackslash}p{0.25\columnwidth}>{\raggedright\arraybackslash}p{0.68\columnwidth}@{}}
\toprule
Class & Operator(s) \\
\midrule
N+1 query & remove \texttt{select\_related}/\texttt{prefetch\_related}, \texttt{prefetch\_related\_objects(...)}, or an eager-loading \texttt{.options(selectinload(...))} \\
Race / concurrency & unwrap a \texttt{with <lock>:} or \texttt{with atomic():} block; drop \texttt{select\_for\_update()}; drop an \texttt{@atomic} decorator \\
Authorisation & delete an \texttt{if <auth-condition>: raise/return} guard or an access-control decorator \\
Unsafe migration & make an added column \texttt{NOT NULL} without a default \\
Resource leak & turn \texttt{with open(...) as f:} (and other closeable resources) into \texttt{f = open(...)}; delete an explicit \texttt{x.close()} \\
Error swallowing & replace a re-raise in an \texttt{except} block by \texttt{pass}; broaden \texttt{except E:} to \texttt{except Exception:} in a non-re-raising handler; delete a validation guard \texttt{if ...: raise} \\
\bottomrule
\end{tabular}
\end{table}

\subsection{Execution-based admission}
For every task we start its published Docker image (the repository at the
pull request's parent commit with dependencies installed), apply the gold patch
and render it again with \texttt{git diff}; the variants are rendered by the
same command, so formatting cannot reveal which patch is the original.  We then
apply the test patch and run the pull request's test files (the files touched by
the test patch and the files of the FAIL\_TO\_PASS tests) with the task's
\texttt{pytest} command.  A task is kept only if at least one test passes with the
gold patch and every FAIL\_TO\_PASS test that appears in the test report passes.  A candidate is \emph{admitted} if every test that passes
with the gold patch also passes with the candidate; otherwise it is recorded as
\emph{caught by tests}.  The execution oracle therefore scores exactly 0\% on
\RJB{} by construction: on the PR's tests, gold and variant are
indistinguishable.

Following SWE-Bench+~\cite{aleithan2024swebenchplus}, we flag \emph{solution
leakage} when a non-trivial line ($\geq$30 characters) added by the gold patch
appears verbatim (up to whitespace) in the issue text, and exclude such tasks
from judging.

\subsection{Repository context retrieval (mitigation, O6)}
To test whether judges fail because they cannot see the code around the diff,
we retrieve context from the \emph{parent} commit, which is neutral with respect
to the two candidates.  The retriever collects (i) definitions of non-builtin
symbols called on changed lines of either candidate and (ii) call sites of
functions modified by either candidate, found with \texttt{git grep} outside
test directories.  Definitions are packed first, then call sites, up to a
character budget $b\in\{2000, 6000\}$ (about 0.6k and 1.7k tokens).  As a
comparison we also provide the parent-commit versions of the touched files,
truncated to $b=6000$ characters (``file context'').

\subsection{Judging protocol}
We use the pairwise prompt distributed with
\CJB{}~\cite{jiang2025codejudgebench} verbatim, with the issue as the user
prompt and the two unified diffs as the assistants' answers.  The prompt asks
for a binary verdict \texttt{[[A>B]]} or \texttt{[[B>A]]}.  Because all judges
run on CPU, we do not let them generate a free-form explanation in the main
condition.  Instead, the assistant turn is forced to begin with
``\texttt{My final verdict is: [[}'' and we read the next-token logits of the
two tokens that start \texttt{A>B]]} and \texttt{B>A]]}.  This yields a greedy
verdict (the larger logit) and a verdict probability
$p_A=\sigma(\ell_A-\ell_B)$.  We check, per tokenizer, that the prefix tokenises
identically before both continuations.  Every pair is judged in both orders
(gold as A, gold as B).  A pair's accuracy is the mean of its two greedy
verdicts, as in \CJB{}.

\textbf{Point-wise condition.}  Each patch is rated alone on a 1--5 scale with
the same forced-prefix technique, and we take the expected rating under the
distribution over the five rating tokens.  A pair is correct if the gold patch
gets the higher expected rating (ties count 0.5).

\textbf{Thinking condition.}  For the hybrid Qwen3 judge we also run its
thinking mode: the model first generates up to \ThinkBudget{} reasoning tokens
(greedy), the reasoning block is closed, and the verdict is read as above.

\subsection{Metrics}
\textbf{Accuracy} is reported with 95\% percentile bootstrap confidence
intervals over pairs~\cite{efron1979bootstrap} (2{,}000 resamples).
\textbf{Position bias} is the accuracy gap between presenting the gold patch
first and second; the \textbf{flip rate} is the share of pairs whose greedy
verdict picks the same \emph{slot} in both orders (and is therefore right in
exactly one order).  \textbf{Self-consistency} is the probability that two
independent temperature-1 samples of the verdict token (restricted to the two
verdict tokens) agree,
$p_A^2+(1-p_A)^2$, averaged over prompts; because the verdict is a single token,
this equals the empirical agreement rate under repeated sampling without having
to sample.  \textbf{Scaling (O3)}: we fit a logistic regression of per-order
correctness on $\log_2$(changed lines), number of files and number of hunks of
the gold patch, with standard errors clustered by pair, and report the diff size
at which the 95\% interval of predicted accuracy first contains 0.5.
